\documentclass[11pt]{article}
\usepackage[margin=0.75in]{geometry}
\usepackage{amsmath,amssymb,amsthm}
\usepackage{array,booktabs,xcolor,colortbl}
\usepackage{tikz}
\usepackage{hyperref}
\setlength{\parindent}{0pt}
\newcommand{\stepnum}[1]{%
  \tikz[baseline=(n.base)] \node[draw,circle,inner sep=1.6pt,
    font=\scriptsize\bfseries,fill=white] (n) {#1};}%
\newtheorem{theorem}{Theorem}
\theoremstyle{definition}
\newtheorem{definition}{Definition}
\title{Flow Proof Artifact: ring\_buffer\_fifo}
\author{Generated by \texttt{flow doc proof}}
\date{\today}
\begin{document}
\maketitle
\bigskip
\noindent{\large \textbf{Derived fact 8.} \textit{Rb\_matches\_queue} \hfill \texttt{rb\_matches\_queue}}\par
\medskip
\noindent\fcolorbox{black!8}{%
\begin{minipage}{\dimexpr\textwidth-2\fboxsep-2\fboxrule}
\textbf{Goal.} We're showing that queue\_order(rb) equals q.items.\\[0.35em]
$\displaystyle \forall rb \in Queue<i32>,\; queue_order(rb) = q.items$
\end{minipage}%
}\par
\medskip
\noindent
\begin{tabular}{@{} c >{\raggedright\arraybackslash}p{0.40\textwidth} @{\hspace{0.4em}\color{gray!50}\vrule width 0.5pt\hspace{0.4em}} c >{\raggedleft\arraybackslash}p{0.40\textwidth} @{}}
\multicolumn{2}{l}{\textbf{Proof}} & \multicolumn{2}{r}{\textbf{Mathematics}} \\
\midrule
\stepnum{1} & We can deduce that ring\_size(rb) equals q.len. & \stepnum{1} & $ring_size(rb) = q.len$ \\[0.55em]
\stepnum{2} & We can deduce that queue\_order(rb) equals q.items. Hence proven. & \stepnum{2} & $queue_order(rb) = q.items$ \\[0.55em]
\bottomrule
\end{tabular}
\label{thm:rb_matches_queue}
\par\medskip\hrule
\bigskip
\noindent{\large \textbf{Derived fact 9.} \textit{Push\_preserves\_fifo} \hfill \texttt{push\_preserves\_fifo}}\par
\medskip
\noindent\fcolorbox{black!8}{%
\begin{minipage}{\dimexpr\textwidth-2\fboxsep-2\fboxrule}
\textbf{Goal.} We're showing that rb\_matches\_queue(rb2, q2).\\[0.35em]
$\displaystyle \forall rb \in \mathbb{Z},\; rb_matches_queue(rb2, q2)$
\end{minipage}%
}\par
\medskip
\noindent
\begin{tabular}{@{} c >{\raggedright\arraybackslash}p{0.40\textwidth} @{\hspace{0.4em}\color{gray!50}\vrule width 0.5pt\hspace{0.4em}} c >{\raggedleft\arraybackslash}p{0.40\textwidth} @{}}
\multicolumn{2}{l}{\textbf{Proof}} & \multicolumn{2}{r}{\textbf{Mathematics}} \\
\midrule
\stepnum{1} & We invoke the derived fact: rb\_matches\_queue (instantiated for rb, q). & & \multicolumn{1}{p{0.40\textwidth}}{} \\[0.55em]
\stepnum{2} & We invoke the derived fact: not ring\_is\_full (instantiated for rb). & & \multicolumn{1}{p{0.40\textwidth}}{} \\[0.55em]
\stepnum{3} & Let rb2 = ring\_push(rb, x). & & \multicolumn{1}{p{0.40\textwidth}}{} \\[0.55em]
\stepnum{4} & Let q2 = queue\_push(q, x). & & \multicolumn{1}{p{0.40\textwidth}}{} \\[0.55em]
\stepnum{5} & From ①, ②, ③, and ④, this implies rb\_matches\_queue(rb2, q2). Hence proven. & \stepnum{5} & $rb_matches_queue(rb2, q2)$ \\[0.55em]
\bottomrule
\end{tabular}
\label{thm:push_preserves_fifo}
\par\medskip\hrule
\end{document}
